\documentclass[10pt,a4paper]{amsart}
\usepackage[margin=19mm]{geometry}
\usepackage{amsmath,amssymb,mathtools}
\usepackage[scr=rsfs,cal=euler]{mathalfa}
\usepackage{xfrac}
\usepackage[unicode,hidelinks,bookmarks=false]{hyperref}
\newcommand{\ie}{i.e.\ }
\newcommand{\cf}{cf.\ }
\newcommand{\resp}{respectively\ }
\newcommand{\Subsection}{\subsection}
\providecommand{\nobreakdash}{\nobreak\hbox{-}}
\setlength{\emergencystretch}{2em}
\multlinegap0pt
\usepackage{amssymb}
\usepackage{enumerate}
\usepackage{float}
\usepackage{xcolor}
\newtheorem{lemma}{Lemma}
\title{Sharp Lower Bounds for Dyadic Square Functions of indicator functions of sets}
\author{Natanael Alpay, Paata Ivanisvili}
\date{Source: arXiv:2502.16045v2 (CC BY 4.0)}
\begin{document}
\maketitle

\noindent\emph{Source and licence.} Sharp Lower Bounds for Dyadic Square Functions of indicator functions of sets. Version arXiv:2502.16045v2 (CC BY 4.0), distributed by arXiv under the Creative Commons Attribution 4.0 International licence, http://creativecommons.org/licenses/by/4.0/, as recorded in the arXiv licence metadata for this exact version and shown on the abstract page https://arxiv.org/abs/2502.16045v2. Full licence notice: \texttt{licenses/arxiv-2502.16045-CC-BY-4.0.txt}.\par\smallskip

\noindent\textit{Editorial scope.} Counted unit: the article's lemma lem:Utilde-sharp-constants-S2 (source0.tex lines 2353--2378). The article's complete original proof of it is delivered with it (source0.tex lines 2380--2614, one \texttt{proof} environment of 235 lines). No mathematics is written, repaired or re-derived: the statement and the proof are the article's own lines, in the article's own notation, and no retained line is substituted or re-flowed.  Retained uncounted as the source-local context the proof needs: lines 2253--2268; lines 2309--2314; lines 2316--2320; lines 2332--2337.  Nothing was omitted from any retained span, so the package carries no omission record.  No retained line contains a citation, so the wrapper carries no bibliography and no bibliography entry.  No retained line contains a figure environment, an image-inclusion command or drawing code, so no image is delivered and the package declares no asset dependency.  Editorial formatting only, in addition: an AMS \texttt{amsart} preamble that restates the article's own declarations and macros used by the retained lines, namely the environments \texttt{lemma} on separate counters, together with the standard packages those lines need.  The retained environments print in this document as Lemma 1 in order of appearance, whereas the article prints the same items with the numbering of its own full document.  The complete native source of the article as distributed in the arXiv e-print for this exact version is bundled unchanged as \texttt{sources/173697/source0.tex}.
\begin{lemma}\label{lem:Utilde-asymptotic-S2}
We have
\begin{equation}
\label{eq:U-exten-sin}
\tilde U(p,0)=\mathbb E_p\sqrt{\tau}
=
\frac{2\sqrt{2\pi}}{\pi^2}\sum_{k=0}^\infty
\frac{\sin((2k+1)\pi p)}{(2k+1)^2},
\end{equation}
and
\[
\mathbb E_p\sqrt{\tau}\asymp p^*\log_2\frac1{p^*},
\qquad p^*:=\min\{p,1-p\}.
\]
In particular, the map $p\mapsto \tilde U(p,0)$ is continuous on $[0,1]$.
\end{lemma}

\begin{equation}
\label{eq:S(p)-rep}
\sum_{k=0}^\infty \frac{\sin((2k+1)x)}{(2k+1)^2}
=
\mathrm{Cl}_2(x)-\frac14\mathrm{Cl}_2(2x),
\end{equation}

\begin{equation}
\label{eq:cl2}
\mathrm{Cl}_2(x):=\sum_{n=1}^\infty \frac{\sin(nx)}{n^2}
=-\int_0^x \log\!\bigl(2\sin(t/2)\bigr)\,dt
\end{equation}

\begin{equation}
\label{eq:S(p)-approx}
\sum_{k=0}^\infty \frac{\sin((2k+1)x)}{(2k+1)^2}
=
\frac{x}{2}\log\frac1x+O(x),
\end{equation}

\begin{lemma}\label{lem:Utilde-sharp-constants-S2}
Let
$
p^*:=\min\{p,1-p\}.
$ Then for every $p\in[0,1]$,
\[
\sqrt{\frac2\pi}\,\ln 2 \; p^*\log_2\frac1{p^*}
\;\le\;
\tilde U(p,0)
\;\le\;
\frac{4\sqrt2\,G}{\pi^{3/2}}\,p^*\log_2\frac1{p^*},
\]
where
$
G=\sum_{k=0}^\infty \frac{(-1)^k}{(2k+1)^2}
$
is Catalan's constant.
Both constants are sharp:
\[
\lim_{p^*\downarrow0}
\frac{\tilde U(p,0)}{p^*\log_2\frac1{p^*}}
=
\sqrt{\frac2\pi}\,\ln 2,
\]
and equality in the upper bound holds at $p=\frac12$.
\end{lemma}

\begin{proof}
By Lemma~\ref{lem:Utilde-asymptotic-S2}, we can write
\[
\tilde U(p,0)=\mathbb E_p\sqrt{\tau}
=
\frac{2\sqrt2}{\pi^{3/2}}\,S(p),
\quad
\text{where}
\quad 
S(p):=\sum_{k=0}^\infty \frac{\sin((2k+1)\pi p)}{(2k+1)^2}.
\]
By equation \eqref{eq:S(p)-rep} using $x=\pi p$, we have 
\[
S(p)=\mathrm{Cl}_2(\pi p)-\frac14 \mathrm{Cl}_2(2\pi p),
\]
where $\mathrm{Cl}_2(x)$ is given by \eqref{eq:cl2}.
Moreover, since 
$
\sin((2k+1)\pi(1-p))
=
\sin((2k+1)\pi p)
$
for every $k\ge 0$, we have
\[
S(1-p)=S(p), \qquad p\in[0,1].
\]
Hence it is enough to work on $[0,1/2]$, where $p^*=p$.

Differentiating the Clausen representation on $(0,1)$ gives
\[
S'(p)
=
-\pi\log\!\Bigl(2\sin\frac{\pi p}{2}\Bigr)
+\frac{\pi}{2}\log\!\bigl(2\sin(\pi p)\bigr).
\]
Using
$
\sin(\pi p)=2\sin\frac{\pi p}{2}\cos\frac{\pi p}{2},
$
we obtain
\[
S'(p)
=
\frac{\pi}{2}\log\!\Bigl(\cot\frac{\pi p}{2}\Bigr).
\]
Since $S(0)=0$, integration yields
\[
S(p)=\int_0^{\pi p}\frac12\log\!\Bigl(\cot\frac t2\Bigr)\,dt.
\]
Differentiating once more, and using
\[
\frac{d}{dx}\log\!\Bigl(\cot\frac x2\Bigr)=-\frac1{\sin x},
\]
we get
\[
S''(p)=-\frac{\pi^2}{2\sin(\pi p)},
\qquad 0<p<1.
\]

We split the proof into three steps.\\

\textbf{Step 1: Lower bound.}
For $0<p\le 1/2$,
$
L(p):=S(p)-\frac{\pi}{2}p\log\frac1p.
$
Then
\[
L''(p)
=
-\frac{\pi^2}{2\sin(\pi p)}+\frac{\pi}{2p}.
\]
Since $\sin(\pi p)\le \pi p$ on $[0,1/2]$, it follows that
\[
-\frac{\pi^2}{2\sin(\pi p)}\le -\frac{\pi}{2p}\implies L''(p)\le 0
\]
 and so
$L$ is concave on $(0,1/2]$.

Taking $ p\downarrow0$ and using \eqref{eq:S(p)-approx} we get
$
L(p)\to 0
$ as $ (p\downarrow0).
$
Also,
\[
L\Bigl(\frac12\Bigr)
=
S\Bigl(\frac12\Bigr)-\frac{\pi}{4}\log2
=
G-\frac{\pi}{4}\log2
>0.
\]
Therefore the continuous extension of $L$ to $[0,1/2]$ is concave, equals $0$ at $0$, and is positive at $1/2$.
Therefore
%A concave function lies above the chord joining two points of its graph, hence
\[
L(p)\ge 0
\qquad \text{for all } p\in\Bigl[0,\frac12\Bigr].
\]
Thus 
% \[
% S(p)\ge \frac{\pi}{2}p\log\frac1p,
% \qquad 0\le p\le \frac12.
% \]
by symmetry,
\[
S(p)\ge \frac{\pi}{2}p^*\log\frac1{p^*},
\qquad 0\le p\le 1.
\]
Multiplying by $\frac{2\sqrt2}{\pi^{3/2}}$ and using
$
\log\frac1{p^*}=(\ln2)\log_2\frac1{p^*},
$
we obtain
\[
\tilde U(p,0)\ge \sqrt{\frac2\pi}\,\ln2 \; p^*\log_2\frac1{p^*}.
\]

\textbf{Step 2: Upper bound.}
Set
$
A:=\frac{2G}{\ln2},
$
and for $0<p\le 1/2$, define
\[
V(p):=S(p)-A\,p\log\frac1p,
\quad 
\text{then}
\quad 
V''(p)
=
-\frac{\pi^2}{2\sin(\pi p)}+\frac{A}{p}.
\]
Now $\sin(\pi p)\ge 2p$ on $[0,1/2]$, so
\[
-\frac{\pi^2}{2\sin(\pi p)}
\ge
-\frac{\pi^2}{4p}
\quad 
\implies
\quad 
V''(p)\ge \frac{A-\pi^2/4}{p}.
\]
Since
$
A=\frac{2G}{\ln2}\approx 2.6434
$ and $
\frac{\pi^2}{4}\approx 2.4674,
$
we have $A>\pi^2/4$, and therefore
\[
V''(p)>0
\qquad\text{for all } p\in\Bigl(0,\frac12\Bigr].
\]
So $V$ is convex on $(0,1/2]$.


Taking $ p\downarrow0$ and using \eqref{eq:S(p)-approx} we get
$
V(p)\to 0
$ as $ (p\downarrow0).
$
Moreover,
\[
V\Bigl(\frac12\Bigr)
=
S\Bigl(\frac12\Bigr)-A\cdot \frac12\log2
=
G-\frac{2G}{\ln2}\cdot \frac12\ln2
=
0.
\]
Thus the continuous extension of $V$ to $[0,1/2]$ is convex and vanishes at both endpoints. Hence 
%A convex function lies below the chord joining two points of its graph; here that chord is the horizontal line $y=0$. Hence
\[
V(p)\le 0
\qquad \text{for all } p\in\Bigl[0,\frac12\Bigr].
\]
Therefore
% \[
% S(p)\le A\,p\log\frac1p,
% \qquad 0\le p\le \frac12.
% \]
% B
by symmetry,
\[
S(p)\le A\,p^*\log\frac1{p^*},
\qquad 0\le p\le 1.
\]
Multiplying by $\frac{2\sqrt2}{\pi^{3/2}}$ and using again
$
\log\frac1{p^*}=(\ln2)\log_2\frac1{p^*},
$
gives
\[
\tilde U(p,0)
\le
\frac{4\sqrt2\,G}{\pi^{3/2}}\,p^*\log_2\frac1{p^*}.
\]

\textbf{Step 3: Sharpness.}
The sharpness of the lower constant follows from the asymptotic already proved in Lemma~\ref{lem:Utilde-asymptotic-S2}:
\[
\tilde U(p,0)
=
\sqrt{\frac2\pi}\,p\log\frac1p+O(p)
\qquad (p\downarrow0).
\]
Dividing by $p\log_2(1/p)=\frac{1}{\ln2}p\log(1/p)$ yields
\[
\lim_{p\downarrow0}
\frac{\tilde U(p,0)}{p\log_2\frac1p}
=
\sqrt{\frac2\pi}\,\ln2.
\]
By symmetry, the same limit holds as $p\uparrow1$, so the lower constant cannot be improved.

For the upper bound, equality holds at $p=\frac12$, because
\[
S\Bigl(\frac12\Bigr)=G
\]
and therefore
\[
\tilde U\Bigl(\frac12,0\Bigr)
=
\frac{2\sqrt2}{\pi^{3/2}}\,G
=
\frac{4\sqrt2\,G}{\pi^{3/2}}\cdot \frac12
=
\frac{4\sqrt2\,G}{\pi^{3/2}}\,
\Bigl(\frac12\Bigr)\log_2 2.
\]
This proves sharpness of the upper constant.
\end{proof}

\end{document}
